%%%%%%%%%%%%%%%%%%%%%%%%%%%%%%%%%%%%%%%%%
% Masters/Doctoral Thesis 
% LaTeX Template
% Version 1.43 (17/5/14)
%
% This template has been downloaded from:
% http://www.LaTeXTemplates.com
%
% Original authors:
% Steven Gunn 
% http://users.ecs.soton.ac.uk/srg/softwaretools/document/templates/
% and
% Sunil Patel
% http://www.sunilpatel.co.uk/thesis-template/
%
% License:
% CC BY-NC-SA 3.0 (http://creativecommons.org/licenses/by-nc-sa/3.0/)
%
% Note:
% Make sure to edit document variables in the Thesis.cls file
%
%%%%%%%%%%%%%%%%%%%%%%%%%%%%%%%%%%%%%%%%%

%----------------------------------------------------------------------------------------
%	PACKAGES AND OTHER DOCUMENT CONFIGURATIONS
%----------------------------------------------------------------------------------------

\documentclass[11pt, oneside]{Thesis} % The default font size and one-sided printing (no margin offsets)

\usepackage{listings} % Wenn man Code Listings benutzen will mit z.B. \begin{lstlisting}[language=Java] ... \end{lstlisting}
\usepackage{color}
\usepackage[usenames,dvipsnames,svgnames,table]{xcolor}
\definecolor{lightgray}{rgb}{.9,.9,.9}
\definecolor{darkgray}{rgb}{.4,.4,.4}
\definecolor{purple}{rgb}{0.65, 0.12, 0.82}

\renewcommand{\lstlistlistingname}{Quellcodeverzeichnis}
\renewcommand{\lstlistingname}{Quellcode}

\lstdefinelanguage{JavaFXScript}{
  keywords={typeof, new, true, false, catch, function, return, null, catch, switch, var, if, in, while, do, else, case, break},
  keywordstyle=\color{blue}\bfseries,
  ndkeywords={class, export, boolean, throw, implements, import, this},
  ndkeywordstyle=\color{darkgray}\bfseries,
  identifierstyle=\color{black},
  sensitive=false,
  comment=[l]{//},
  morecomment=[s]{/*}{*/},
  commentstyle=\color{purple}\ttfamily,
  stringstyle=\color{red}\ttfamily,
  morestring=[b]',
  morestring=[b]"
}

\lstdefinelanguage{Java}{
  keywords={package, import, public, private, class, extends, implements, void, throws, new, null, int, double, long, String, this, static, interface, abstract},
  keywordstyle=\color{blue}\bfseries,
  identifierstyle=\color{black},
  sensitive=false,
  comment=[l]{//},
  morecomment=[s]{/*}{*/},
  commentstyle=\color{gray}\ttfamily,
  stringstyle=\color{Green}\ttfamily,
  morestring=[b]',
  morestring=[b]"
}

\usepackage{bibgerm} % Wenn man ein Literaturverzeichnis nach der Deutschen Norm haben will
%\usepackage[square, comma]{natbib} % Use the natbib reference package - read up on this to edit the reference style; if you want text (e.g. Smith et al., 2012) for the in-text references (instead of numbers), remove 'numbers' 
\hypersetup{urlcolor=blue, colorlinks=true} % Colors hyperlinks in blue - change to black if annoying

\newcommand{\thesisTitle}{Dokumententitel}
\newcommand{\authorName}{Student}

% PDF meta-data
\hypersetup{pdftitle=\thesisTitle}
\hypersetup{pdfauthor=\authorName}

\begin{document}

\frontmatter % Use roman page numbering style (i, ii, iii, iv...) for the pre-content pages

\setstretch{1.5} % Line spacing of 1.5

% Define the page headers using the FancyHdr package and set up for one-sided printing
\fancyhead{} % Clears all page headers and footers
\rhead{\thepage} % Sets the right side header to show the page number
\lhead{} % Clears the left side page header

\pagestyle{fancy} % Finally, use the "fancy" page style to implement the FancyHdr headers

\newcommand{\HRULE}{\rule{\linewidth}{0.5mm}} % New command to make the lines in the title page
\newcommand{\HRule}{\rule{\linewidth}{0.1mm}} % New command to make the lines in the title page

%----------------------------------------------------------------------------------------
%	TITLE PAGE
%----------------------------------------------------------------------------------------

%----------------------------------------------------------------------------------------
%	TITLE PAGE
%----------------------------------------------------------------------------------------

\begin{titlepage}

%top
\begin{center}

\includegraphics[width=0.3\textwidth]{images/TH_Koeln_Logo.png}\\

\vspace*{1cm}

\begin{Huge}
Praxisprojekt\\
\end{Huge}

\vspace*{0.5cm}

\begin{LARGE}
Ermittlung relevanter Themengebiete für die Entwicklung eines Tools zur Unterstützung beim Erstellen von Gestaltungslösungen im Hochschulkontext
\end{LARGE}

\vspace*{1.5cm}

\begin{large}
	von\\
    Christian Alexander Poplawski - 11088931\\

\vspace{1.0cm}

	an der Technology, Arts and Sciences TH Köln\\
    Campus Gummersbach\\
    im Studiengang Medieninformatik (Bachelor)\\

\vspace{1.0cm}

    \begin{tabular}{rl}
        Betreuer:  &  Prof. Dipl. Des. Christian Noss\\
       			   &  \small Technology, Arts and Sciences TH Köln \\[1.0em]

	\end{tabular}

\end{large}

\vspace{1.5cm}

Technology, Arts and Sciences TH Köln, \today \\


\vfill
\end{center}

\end{titlepage}


%----------------------------------------------------------------------------------------
%	ABSTRACT PAGE
%----------------------------------------------------------------------------------------

%----------------------------------------------------------------------------------------
%	ABSTRACT PAGE
%----------------------------------------------------------------------------------------
\thispagestyle{plain}

\addtotoc{Abstract} % Add the "Zusammenfassung" page entry to the Contents

\begin{huge}
	\textbf{Abstract}\\
\end{huge}

Hier steht später das Abstract.

\clearpage

%----------------------------------------------------------------------------------------
%	LIST OF CONTENTS
%----------------------------------------------------------------------------------------

\pagestyle{fancy} % The page style headers have been "empty" all this time, now use the "fancy" headers as defined before to bring them back
\hypersetup{linkcolor=black}
\lhead{\emph{Inhaltsverzeichnis}} % Set the left side page header to "Inhaltsverzeichnis"

\tableofcontents % Write out the Table of Contents

\clearpage

%----------------------------------------------------------------------------------------
%	LIST OF FIGURES/TABLES PAGES
%----------------------------------------------------------------------------------------

%----------------------------------------------------------------------------------------
%	LIST OF FIGURES/TABLES PAGES
%----------------------------------------------------------------------------------------

\pagestyle{fancy} % The page style headers have been "empty" all this time, now use the "fancy" headers as defined before to bring them back

\lhead{\emph{Abbildungsverzeichnis}} % Set the left side page header to "List of Figures"
\addtotoc{Abbildungsverzeichnis} % Add the "Abbildungsverzeichnis" page entry to the Contents
\listoffigures % Write out the List of Figures
\addtocontents{toc}{\vspace{-0.5cm}} % Add a gap in the Contents, for aesthetics

\lhead{\emph{Quellcodeverzeichnis}} % Set the left side page header to "List of Tables"
\addtotoc{Quellcodeverzeichnis}
\lstlistoflistings 
\addtocontents{toc}{\vspace{-0.5cm}}

\clearpage

%\lhead{\emph{Tabellenverzeichnis}} % Set the left side page header to "List of Tables"
%\addtotoc{Tabellenverzeichnis} % Add the "Tabellenverzeichnis" page entry to the Contents

%\addtocontents{toc}{\vspace{0.5cm}} % Add a gap in the Contents, for aesthetics
%\addtocontents{toc}{\vspace{-0.5cm}} % Add a gap in the Contents, for aesthetics
%\listoftables % Write out the List of Tables


%----------------------------------------------------------------------------------------
%	ABBREVIATIONS
%----------------------------------------------------------------------------------------

%----------------------------------------------------------------------------------------
%	ABBREVIATIONS
%----------------------------------------------------------------------------------------

\clearpage % Start a new page

\thispagestyle{plain}

\setstretch{1.5} % Set the line spacing to 1.5, this makes the following tables easier to read

\lhead{} % Set the left side page header to ""

\addtotoc{Abkürzungsverzeichnis} % Add the "Abkürzungsverzeichnis" page entry to the Contents
\begin{huge}
	\textbf{Abkürzungsverzeichnis}\\
\end{huge}

%\lhead{\emph{Abkürzungsverzeichnis}} % Set the left side page header to "Abkürzungsverzeichnis"

\begin{tabular}{ l l }
    \textbf{AWT} & Abstract Window Toolkit\\
    \textbf{API} & Application Programming Interface\\
    \textbf{EDT} & Event Dispatch Thread\\
    \textbf{F3} & Form Follows Function\\
    \textbf{HLS} & HTTP Live Streaming\\ 
    \textbf{RIA} & Rich Internet Application\\    
    \textbf{UI} & User Interface\\    
\end{tabular}


%----------------------------------------------------------------------------------------
%	THESIS CONTENT - CHAPTERS
%----------------------------------------------------------------------------------------

\mainmatter % Begin numeric (1,2,3...) page numbering


\pagestyle{fancy} % Return the page headers back to the "fancy" style

% Include the chapters of the thesis as separate files from the Chapters folder
% Uncomment the lines as you write the chapters

% Chapter 1

%\footnote{vgl. \cite{studie}}

\newcommand{\chaptertitle}{Einleitung}

\chapter{\chaptertitle} % Main chapter title

\label{Einleitung} % For referencing the chapter elsewhere, use \ref{Chapter1} 

\lhead{\chaptername{} \thechapter{} - \emph{\chaptertitle}} % This is for the header on each page - perhaps a shortened title

%----------------------------------------------------------------------------------------

\section{Einleitung}

%----------------------------------------------------------------------------------------

\section{Aktueller Forschungsstand}

%----------------------------------------------------------------------------------------

\section{Themenrelevanz}

%----------------------------------------------------------------------------------------

\section{Aufbau und Methodik}

%----------------------------------------------------------------------------------------

\clearpage
% Chapter 4

%\footnote{vgl. \cite{studie}}

\renewcommand{\chaptertitle}{Neues Kapitel}

\chapter{\chaptertitle} % Main chapter title

\label{Kapitel 2} % For referencing the chapter elsewhere, use \ref{Chapter1} 

\lhead{\chaptername{} \thechapter{} - \emph{\chaptertitle}} % This is for the header on each page - perhaps a shortened title

%----------------------------------------------------------------------------------------

\section{Unterkapitel}

\subsection{Unterunterkapitel}

%----------------------------------------------------------------------------------------

\clearpage
% Chapter 4

%\footnote{vgl. \cite{studie}}

\renewcommand{\chaptertitle}{Nächstes Kapitel}

\chapter{\chaptertitle} % Main chapter title

\label{Kapitel 3} % For referencing the chapter elsewhere, use \ref{Chapter1} 

\lhead{\chaptername{} \thechapter{} - \emph{\chaptertitle}} % This is for the header on each page - perhaps a shortened title

%----------------------------------------------------------------------------------------

\section{Unterkapitel}

\subsection{Unterunterkapitel}

%----------------------------------------------------------------------------------------

\clearpage
% Chapter 5

\chapter{Fazit} % Main chapter title

\label{Fazit} % For referencing the chapter elsewhere, use \ref{Aufgabenstellung} 

\lhead{\chaptername{} \thechapter{} - \emph{Fazit}} % This is for the header on each page - perhaps a shortened title

Abschließend lässt sich die Zielsetzung als erreicht ansehen. \\
Es wurden relevante Themengebiete für die Konzeption eines Tools definiert. In diesen Themengebieten wurden Regeln festgehalten, die auch für eine spätere programmatische Umsetzung verwendet werden können. In machen Gebieten, wie der Typographie oder den Farben, lassen diese Regeln mehr Spielraum für Interaktivität zu, in Anderen, wie beispielsweise den Interaktiven Elementen und Bildern, weniger. Lediglich im Bereich Whitespace konnten keinerlei Regeln definiert werden.

Mögliche Ideen für eine Umsetzung wurden in Form von Wireframes festgehalten und im Bereich der Typographie konnte mit Hilfe eines \textit{Proof of Concepts} gezeigt werden, dass auch eine tatsächliche technische Umsetzung mit moderatem Aufwand durchaus möglich ist.

Während des Projektes konnte aber auch festgestellt werden, dass das Übertragen von eher subjektiven Regeln für die Gestaltung von Artefakten in eine objektive, maschinenlesbare Form nicht einfach ist. Viele Entscheidungen, die Designer jeden Tag treffen, kann ein Computer nicht nachvollziehen. Auch können getroffene Entscheidungen nicht einfach als richtige oder falsche Referenzwerte übergeben werden, da deren Richtigkeit stark situationsabhängig ist. \\
Hier stellt sich die Frage, ob der Bereich \textit{Machine Learning} eine sinnvolle Ergänzung für das Projekt bieten könnte, also ob ein Computer, ähnlich wie ein Mensch, ein \textit{Gefühl} für eine gute oder schlechte Gestaltung entwickeln kann, wenn seine Datengrundlage groß genug ist.

Weiterhin ist auch die Frage nach dem tatsächlichen Nutzen den Tools noch ungeklärt. Zwar wird deutlich, das ein solches Tool eine gewisse Relevanz hat und auch umgesetzt werden kann, es kann jedoch nicht Festgestellt werden, ob das Tool dem Nutzer einen tatsächlichen Mehrwert liefern würde. Hierfür wären Nutzertests mit Prototypen nötig. 


\clearpage

%----------------------------------------------------------------------------------------
%	THESIS CONTENT - APPENDICES
%----------------------------------------------------------------------------------------

%----------------------------------------------------------------------------------------
%	THESIS CONTENT - APPENDICES
%----------------------------------------------------------------------------------------

\addtocontents{toc}{\vspace{2em}} % Add a gap in the Contents, for aesthetics

\appendix % Cue to tell LaTeX that the following 'chapters' are Appendices

% Include the appendices of the thesis as separate files from the Appendices folder
% Uncomment the lines as you write the Appendices

%% Appendix A

\chapter{Anforderungen} % Main appendix title

\label{anhangA} % For referencing this appendix elsewhere, use \ref{AppendixA}

\lhead{Anhang A. \emph{Anforderungen}} % This is for the header on each page - perhaps a shortened title

\section{Ursprüngliche Anforderungen}
\label{uanforderungen}

%\input{Appendices/AppendixB}
%\input{Appendices/AppendixC}

\addtocontents{toc}{\vspace{2em}} % Add a gap in the Contents, for aesthetics

\backmatter

\clearpage


%----------------------------------------------------------------------------------------
%	BIBLIOGRAPHY
%----------------------------------------------------------------------------------------

%----------------------------------------------------------------------------------------
%	BIBLIOGRAPHY
%----------------------------------------------------------------------------------------

\label{Literaturverzeichnis}

\bibliographystyle{plainnat} % Use the "unsrtnat" BibTeX style for formatting the Bibliography

\bibliography{Bibliography} % The references (bibliography) information are stored in the file named "Bibliography.bib"

\lhead{\emph{Literaturverzeichnis}} % Change the page header to say "Bibliography"

\clearpage

%----------------------------------------------------------------------------------------
%	DECLARATION PAGE
%	Your institution may give you a different text to place here
%----------------------------------------------------------------------------------------

%%----------------------------------------------------------------------------------------
%	DECLARATION PAGE
%	Your institution may give you a different text to place here
%----------------------------------------------------------------------------------------

\Declaration{

\addtocontents{toc}{\vspace{1em}} % Add a gap in the Contents, for aesthetics

\vspace*{2cm}

Ich versichere, die von mir vorgelegte Arbeit selbständig verfasst zu haben. Alle Stellen, die wörtlich oder sinngemäß aus veröffentlichten oder nicht veröffentlichten Arbeiten anderer entnommen sind, habe ich als entnommen kenntlich gemacht. Sämtliche Quellen und Hilfsmittel die ich für die Arbeit benutzt habe, sind angegeben. Die Arbeit hat mit gleichem Inhalt bzw. in wesentlichen Teilen noch keiner anderen Prüfungsbehörde vorgelegen.\\

\vspace*{2cm}
 
\rule[1em]{25em}{0.5pt}\\ % This prints a line for the signature

\vspace*{-1.2cm}

Datum, Ort\\

\vspace*{1.5cm}

\rule[1em]{25em}{0.5pt}\\ % This prints a line to write the date

\vspace*{-1.2cm}

Unterschrift\\
}

\clearpage % Start a new page

\end{document}